% "высота букв, цифр и других знаков - не менее 12пт"
\documentclass[a4paper,14pt]{extarticle}

% Кодировка и язык
\usepackage{polyglossia}
\setdefaultlanguage{russian}
\setotherlanguage{english}

% Шрифты
\usepackage{fontspec}
\defaultfontfeatures{Ligatures=TeX}
\setmainfont{Times New Roman}
\newfontfamily\cyrillicfont{Times New Roman}[Ligatures=TeX]
\newfontfamily\cyrillicfonttt{Times New Roman}[Ligatures=TeX]

% Настройка полей страницы
\usepackage{geometry}
\geometry{
	left=30mm,
	right=15mm,
	top=20mm,
	bottom=20mm
}

% Межстрочный интервал 1.5
\usepackage{setspace}
\onehalfspacing

% Абзацный отступ
\usepackage{indentfirst}
\setlength{\parindent}{1.25cm}

% Нумерация страниц
\usepackage{fancyhdr}
\pagestyle{fancy}
\fancyhf{}
\fancyfoot[C]{\thepage}
\renewcommand{\headrulewidth}{0pt}
\renewcommand{\footrulewidth}{0pt}

% Ограничение глубины содержания (только 1.1, без 1.1.1)
\setcounter{tocdepth}{2}

% Настройка заголовков (ГОСТ)
\usepackage{titlesec}

% Настраиваем абзацный отступ для всех уровней
\newlength{\gostindent}
\setlength{\gostindent}{1.25cm}

% --- РАЗДЕЛЫ ---
% Нумерованные: с прописной буквы, жирный, 14пт, абзацный отступ
\titleformat{\section}
{\normalfont\fontsize{14pt}{\baselineskip}\selectfont\bfseries}
{\thesection}
{0.5em}
{}
\titlespacing*{\section}{\gostindent}{0pt}{0pt}

% НЕНУМЕРОВАННЫЕ разделы (Введение, Список лит-ры) — ПО ЦЕНТРУ
\titleformat{name=\section,numberless}
{\normalfont\fontsize{14pt}{\baselineskip}\selectfont\bfseries\centering}
{}
{0pt}
{\MakeUppercase}
\titlespacing*{name=\section,numberless}{0pt}{0pt}{0pt}

% --- ПОДРАЗДЕЛЫ ---
\titleformat{\subsection}
{\normalfont\fontsize{14pt}{\baselineskip}\selectfont\bfseries}
{\thesubsection}
{0.5em}
{}
\titlespacing*{\subsection}{\gostindent}{1em}{0.5em}

% --- ПОДПУНКТЫ ---
\titleformat{\subsubsection}
{\normalfont\fontsize{14pt}{\baselineskip}\selectfont}
{\thesubsubsection}
{0.5em}
{}
\titlespacing*{\subsubsection}{\gostindent}{0pt}{0pt}

% Каждый раздел с новой страницы
\newcommand{\sectionbreak}{\clearpage}

% Графика и рисунки
\usepackage{graphicx}
\usepackage{caption}
\usepackage{subcaption}

% Настройка подписей к рисункам
\DeclareCaptionLabelSeparator{emdash}{ --- }
\captionsetup[figure]{
	labelsep=emdash,
	justification=centering,
	font=normalsize,
	skip=6pt  % Отступ между рисунком и подписью
}

% Настройка подписей к таблицам
\captionsetup[table]{
	labelsep=emdash,
	justification=raggedright,
	singlelinecheck=false,
	font=normalsize,
	skip=6pt  % Отступ между подписью и таблицей
}

% Таблицы
\usepackage{array}
\usepackage{longtable}
% Настройка подписи для longtable по ГОСТ
\setlength\LTleft{0pt}
\setlength\LTright{\fill}
\captionsetup[longtable]{
	labelsep=emdash,
	justification=raggedright,
	singlelinecheck=false,
	font=normalsize,
	skip=6pt
}
\usepackage{multirow}

% Математика
\usepackage{amsmath}
\usepackage{amssymb}

% Списки: БЕЗ ТОЧЕК, только тире или скобки
\usepackage{enumitem}
% Общая настройка для всех списков
\setlist{
	nosep,                      % Убираем лишние вертикальные отступы
	labelindent=\parindent,     % Позиция номера = абзацный отступ (1.25см)
	leftmargin=*,               % Текст после первой строки идет от левого поля
	align=left                  % Выравнивание номера по левому краю отступа
}

% Настройка названия рисунков
\addto\captionsrussian{\renewcommand{\figurename}{Рисунок}}

% Нумерованные списки: «1. текст», номер с абзацного отступа,
% между номером и текстом ровно 0.5em (около одного символа в TNR 14pt).
% leftmargin=* — enumitem сам считает ширину блока, текст продолжается ровно
% под началом текста первого item'а.
\setlist[enumerate]{
	label=\arabic*.,
	ref=\arabic*,
	labelindent=\parindent,
	labelsep=0.5em,
	leftmargin=*,
	align=left
}

% Маркированные списки (тире)
\setlist[itemize]{
	label=---
}

% Гиперссылки (для кликабельного содержания)
\usepackage{hyperref}
\hypersetup{
	colorlinks=true,
	linkcolor=black,
	citecolor=black,
	urlcolor=blue,
	breaklinks=true
}
\usepackage{xurl}
\sloppy

% Для вставки PDF в приложение
\usepackage{pdfpages}

% ------------
%  содержание
% ------------
\usepackage{tocloft}
\usepackage{setspace}

% Вертикальный воздух между пунктами TOC.
% Внутри пункта строки идут одинарным интервалом (см. \fontsize ниже),
% а между пунктами добавляется небольшой пробел, имитирующий полуторный.
\setlength{\cftbeforesecskip}{0pt}
\setlength{\cftbeforesubsecskip}{0pt}
\setlength{\cftbeforesubsubsecskip}{0pt}

% отступы:
% > разделы 0
% > подразделы 0.49см
% > подпункты 0.99см
\setlength{\cftsecindent}{0pt}
\setlength{\cftsubsecindent}{0.49cm}
\setlength{\cftsubsubsecindent}{0.99cm}

% Ширина колонки номера
\setlength{\cftsecnumwidth}{0.75em}
\setlength{\cftsubsecnumwidth}{1.55em}
\setlength{\cftsubsubsecnumwidth}{2.7em}

% Плотность отточия
\renewcommand{\cftdotsep}{1.1}

% Ширина колонки номера страницы (2 цифры в кегле 14pt помещаются в 1.3em)
\cftsetpnumwidth{0.85em}
\cftsetrmarg{0.82em}

\renewcommand{\cfttoctitlefont}{\hfill\fontsize{14pt}{\baselineskip}\selectfont\bfseries}
\renewcommand{\cftaftertoctitle}{\hfill}
\renewcommand{\cftaftertoctitleskip}{1em}

% Отточие идёт вплотную к номеру страницы (без принудительного \hspace)
\renewcommand{\cftsecleader}{\cftdotfill{\cftdotsep}}
\renewcommand{\cftsubsecleader}{\cftdotfill{\cftdotsep}}
\renewcommand{\cftsubsubsecleader}{\cftdotfill{\cftdotsep}}

\addto\captionsrussian{%
	\renewcommand{\contentsname}{СОДЕРЖАНИЕ}%
}

% Шрифт записей TOC: одинарный интервал (baselineskip = 17pt при кегле 14pt).
% Так многострочные пункты не разъезжаются.
\renewcommand{\cftsecfont}{\normalfont\fontsize{14pt}{14pt}\selectfont}
\renewcommand{\cftsubsecfont}{\normalfont\fontsize{14pt}{14pt}\selectfont}
\renewcommand{\cftsubsubsecfont}{\normalfont\fontsize{14pt}{14pt}\selectfont}

% Шрифт номеров страниц: тот же одинарный интервал, без принудительной ширины
\renewcommand{\cftsecpagefont}{\normalfont\fontsize{14pt}{14pt}\selectfont}
\renewcommand{\cftsubsecpagefont}{\normalfont\fontsize{14pt}{14pt}\selectfont}
\renewcommand{\cftsubsubsecpagefont}{\normalfont\fontsize{14pt}{14pt}\selectfont}

% Гасим растягивание пробелов в многострочных пунктах TOC.
% tocloft переопределяет \l@section/\l@subsection/\l@subsubsection своими
% макросами, которые ставят \rightskip \@tocrmarg (фиксированное правое поле)
% и тем самым justify-ят все строки заголовка. Через etoolbox\patchcmd
% заменяем правое поле на упругое \@tocrmarg plus 1fil — теперь нерастянутые
% строки оставляют пустое пространство справа, без растягивания пробелов.
\usepackage{etoolbox}
\makeatletter
\patchcmd{\l@section}
  {\rightskip \@tocrmarg}{\rightskip \@tocrmarg plus 1fil}{}{}
\patchcmd{\l@subsection}
  {\rightskip \@tocrmarg}{\rightskip \@tocrmarg plus 1fil}{}{}
\patchcmd{\l@subsubsection}
  {\rightskip \@tocrmarg}{\rightskip \@tocrmarg plus 1fil}{}{}
\makeatother

% Команда для вставки приложения в текст и TOC
\newcommand{\appsection}[2]{% #1 - Буква (А, Б..), #2 - Название
	\newpage
	\begin{center}
		\textbf{Приложение #1}
	\end{center}
	\begin{center}
		\textbf{#2}
	\end{center}

	% Временно отключаем отточие (лидеры) для этой записи в TOC
	\addtocontents{toc}{\protect\renewcommand{\protect\cftsecleader}{\protect\hfill}}

	% Добавляем в содержание формат: ПРИЛОЖЕНИЕ Х (Название)
	\addcontentsline{toc}{section}{ПРИЛОЖЕНИЕ #1 (#2)}

	% Возвращаем отточие обратно для последующих записей (если будут)
	\addtocontents{toc}{\protect\renewcommand{\protect\cftsecleader}{\protect\cftdotfill{\protect\cftdotsep}}}
}

\usepackage{float}

% -------------------
%    bibliography
% -------------------
\usepackage[backend=biber,style=gost-numeric,sorting=none]{biblatex}
\usepackage{etoolbox}

% Запрет переноса между инициалами и фамилией ("А. Ф.~Стефаниди"),
% а также между двумя инициалами. Норма русской научной верстки.
\setcounter{highnamepenalty}{10000}
\setcounter{lownamepenalty}{10000}

\AtBeginBibliography{\hypersetup{hidelinks}}
\DeclareFieldFormat{author}{\textnormal{#1}}
\DeclareFieldFormat{editor}{\textnormal{#1}}
\renewbibmacro*{author}{%
  \printnames[author]{author}%
}

% ГОСТ 7.32 п.4.10.3: нумеровать арабскими цифрами с точкой и печатать с
% абзацного отступа. Номер «N.» стоит с абзацного отступа (1.25 см) —
% как обычный абзац, — текст идёт сразу за номером, переносы на новую
% строку выравниваются по левому полю (висячий отступ).
\defbibenvironment{bibliography}
  {\list
     {\printfield{labelnumber}.}              % метка вида «1.»
     {\setlength{\labelwidth}{0pt}%
      \setlength{\labelsep}{0.4em}%           % один пробел между «1.» и текстом
      \setlength{\leftmargin}{0pt}%           % перенос — по левому полю
      \setlength{\itemindent}{1.25cm}%        % первая строка — с абзацного отступа
      \setlength{\itemsep}{0pt}%
      \setlength{\parsep}{0pt}%
      \setlength{\topsep}{0pt}%
      \setlength{\partopsep}{0pt}%
      \setlength{\listparindent}{1.25cm}%
      \setlength{\rightmargin}{0pt}%
      \setlength{\parskip}{0pt}%
      \setlength{\parindent}{0pt}%
      \renewcommand{\makelabel}[1]{##1\hfil}% выравнивание метки по левому краю
     }}
  {\endlist}
  {\item}

\renewrobustcmd*{\mkbibemph}[1]{\normalfont#1}
\ExecuteBibliographyOptions{maxbibnames=99, minbibnames=1}
\addbibresource{lib/biblio.bib}
\renewcommand*{\bibfont}{\setlength{\parindent}{1.25cm}%
\setmainfont{Times New Roman}\fontsize{14pt}{16pt}\selectfont}
\usepackage{booktabs}
\usepackage{pdflscape}
\DefineBibliographyStrings{russian}{
    urlseen = {дата обращения} % Текст перед датой обращения
}

% Переопределяем, как выводится URL и дата обращения
% Это более сложная настройка, если стандартный вывод gost-numeric не устраивает.
% Сначала попробуйте просто с urldate и note/[howpublished] в .bib файле.
% Если нужно тонко настроить, то можно так:
\DeclareFieldFormat[online]{url}{%
  {[Электронный ресурс]. URL:\space\nolinkurl{#1}}% <--- \nolinkurl убирает кликабельность
}
% или если всегда добавлять [Электронный ресурс] перед URL для @online

\DeclareFieldFormat{urldate}{\mkbibparens{\bibstring{urlseen}\addcolon\space#1}}

% Чтобы "URL:" появлялось перед самой ссылкой
% (стили ГОСТ иногда требуют это, иногда нет)
\DefineBibliographyStrings{russian}{
  url = {URL} % Меняем "Режим доступа" на "URL"
}

\DeclareFieldFormat[article,inproceedings,incollection,book,thesis,report,patent,online,misc,unpublished]{title}{\normalfont{#1}} % Явно для основных типов
\DeclareFieldFormat{booktitle}{\normalfont{#1}}    % Для названия сборника или трудов конференции
\DeclareFieldFormat{journaltitle}{\normalfont{#1}} % Для названия журнала
\DeclareFieldFormat{issuetitle}{\normalfont{#1}}   % Для названия выпуска журнала
\DeclareFieldFormat{eventtitle}{\normalfont{#1}}   % Для названия мероприятия (конференции)
\setlength{\bibindent}{1.25cm} % Отступ слева для каждого элемента списка

